% GENERATED, edit the French unit or the translation table
\mxnewpage\section*{EV1-B -- Common test}\label{EV1-B}
\addcontentsline{toc}{section}{EV1-B -- Common test}
\textbf{Paper reserved for the teacher}. Chapters 1 to 4.\par Duration: 2 h. No calculator. Marked out of 20. 14 November 2026.
\subsection*{Part 1 -- Fluency (6 points)}
Twelve questions, 0.5 point each. Only one correct answer; no justification required. An incorrect or missing answer scores 0, with no loss of points.
\begin{enumerate}
\item Calculate $6+2\times3$.\par A. $11$\quad B. $12$\quad C. $15$\quad D. $24$
\item Calculate $\frac12+\frac13$.\par A. $\frac{1}{6}$\quad B. $\frac{5}{6}$\quad C. $\frac{2}{5}$\quad D. $\frac{5}{3}$
\item Calculate $\dfrac{2^6}{2^{4}}$.\par A. $2$\quad B. $11$\quad C. $1$\quad D. $4$
\item A value of 50 increases by 20 \%. Final value?\par A. $60$\quad B. $70$\quad C. $10$\quad D. $40$
\item A value of 80 is obtained after a 20 \% decrease. Initial value?\par A. $64$\quad B. $96$\quad C. $60$\quad D. $100$
\item A 25 \% increase then a 20 \% decrease: overall coefficient?\par A. $\frac{21}{20}$\quad B. $\frac{4}{5}$\quad C. $1$\quad D. $\frac{9}{10}$
\item The solution of $2x+6=0$ is\par A. $3$\quad B. $12$\quad C. $-3$\quad D. $-12$
\item The slope between $(0;6)$ and $(2;12)$ is\par A. $-3$\quad B. $3$\quad C. $6$\quad D. $2$
\item The mean of $6, 8, 10$ is\par A. $8$\quad B. $24$\quad C. $6$\quad D. $9$
\item A sector of $90^\circ$ represents what fraction of the total?\par A. $\frac{1}{3}$\quad B. $\frac{1}{2}$\quad C. $\frac{1}{4}$\quad D. $\frac{3}{4}$
\item A probability is $\dfrac{2}{5}$. Its complement is\par A. $\frac{2}{5}$\quad B. $\frac{3}{5}$\quad C. $\frac{5}{2}$\quad D. $\frac{1}{5}$
\item For $f(x)=x^2+6$, $f(-2)$ is\par A. $8$\quad B. $24$\quad C. $10$\quad D. $2$
\end{enumerate}
\mxnewpage\subsection*{Part 2 -- Independent exercises (14 points)}
Every step must be justified. Exact results are expected unless stated otherwise.
\mxneed{12}
\subsection*{Exercise 1 -- Growth by pattern (5 points)}
A family of patterns is indexed by the integers $n\geq1$. The number of dots in the pattern of index $n$ is given by $u_n=n^2+2n$. We are interested in particular in the first pattern with at least 99 dots; no figure is needed to use the formula given.
\begin{enumerate}
\item Calculate $u_1,u_2,u_3$. \hfill(1 pt)
\item Calculate $u_{n+1}-u_n$ then establish the variation. \hfill(2 pts)
\item From which index does the pattern contain at least 99 dots? Justify without an endless sweep. \hfill(2 pts)
\end{enumerate}
\mxneed{12}
\subsection*{Exercise 2 -- Two dimensions and a constraint (5 points)}
A rectangle has a perimeter of 26 m and an area of 40 square metres. Let $x$ and $y$ be its two dimensions, in metres. Both numbers are strictly positive. We will then compare this rectangle with the other rectangles having the same perimeter.
\begin{enumerate}
\item Write the relations between the two lengths x and y. \hfill(1 pt)
\item Determine the dimensions using a quadratic equation. \hfill(2 pts)
\item Among the rectangles with the same perimeter, determine the maximum area and compare it with 40. \hfill(2 pts)
\end{enumerate}
\mxneed{12}
\subsection*{Exercise 3 -- Angle then approximation (4 points)}
Questions 1 and 2 concern two vectors $\vec u$ and $\vec v$ with norms 2 and 3 respectively, such that $\vec u\cdot\vec v=3$. Question 3 is independent: it concerns approximation by a tangent to the square function.
\begin{enumerate}
\item Determine the measure in degrees of the unoriented angle between $\vec u$ and $\vec v$. \hfill(1 pt)
\item Calculate the exact value of $\norm{\vec u-\vec v}$. \hfill(1 pt)
\item Independently, the tangent to the square function at 1 has equation $y=2x-1$. Use it to approximate $1.03^2$, then calculate the difference between the exact value and the approximation. \hfill(2 pts)
\end{enumerate}
